% Post-CP22 candidate. No assertion is imported from characteristic-zero sources.
\subsection{Three total ramifications force ordinarity in genus two}
\status{Proved and independently checked.}

\begin{lemma}[Three totally ramified fibres of a trigonal map]
\label{lem:char2-genus-two-three-total-ramifications-candidate}
\textup{[Proved under the stated hypotheses; complete independent audits passed.]}
Let $C$ be a smooth projective curve of genus two over an algebraically
closed field of characteristic two. If a degree-three morphism
$C\to\mathbf P^1$ has at least three distinct totally ramified
points, then $C$ is ordinary.
\end{lemma}
\begin{proof}
The morphism is separable, since its degree is odd. The different has
degree eight by Riemann--Hurwitz. Each of the stated tame ramification
points has different exponent two. The remaining different degree is
two. Consequently there are exactly two possibilities: a fourth tame
point of ramification index three, or one point of ramification index
two and different exponent two. A wild quadratic point has different
exponent at least two, so there are no other possibilities.
These different formulas are the characteristic-independent statements
in \citet[Section~53.12, especially Lemmas~53.12.2 and~53.12.4]{Stacks0C1B}.

Let $Y$ be the smooth Galois closure. Its transitive permutation group
is either $C_3$ or $S_3$. If all four ramification points are tame of
index three, the sign quotient has no ramification. A nontrivial
unramified double cover of $\mathbf P^1$ is impossible by
Riemann--Hurwitz. Thus the group is $C_3$ in this case.
Choose infinity away from the four branch points. Kummer theory
\citep[Lemma~9.24.1, Tag~09DX]{Stacks09DX}, which applies because
three is prime to the characteristic, gives
\[
                   y^3=F(z)/G(z),
\]
where $F$ and $G$ are monic separable quadratics with no common root.
Indeed the four nonzero branch exponents are two $+1$ and two $-1$
modulo three, since their sum is zero. Put
$\eta=yG$, so that $\eta^3=FG^2$.
The two forms
\[
             \omega=dz/\eta,\qquad \omega'=G\,dz/\eta^2
\]
are regular and independent. At a branch point their orders are zero
or one; at infinity both have order zero. They therefore form a
basis of the genus-two space of regular differentials.
Write $F=z^2+f_1z+f_0$ and $G=z^2+g_1z+g_0$.
The Cartier rules in characteristic two give
\[
              C\omega=\sqrt{f_1}\,\omega',\qquad
              C\omega'=\sqrt{g_1}\,\omega.
\]
For example $C(dz/\eta)=\eta^{-2}C(FG^2\,dz)
=\eta^{-2}G C(F\,dz)$. Both $f_1$ and $g_1$ are nonzero,
because the respective quadratics are separable. Cartier is
therefore bijective, proving ordinarity in the cyclic case.

In the remaining case the Galois group is $S_3$. Let
$Z=Y/C_3$ be its quadratic sign quotient. It is unramified above the
three tame branch values. At the remaining branch value its completed
extension is exactly the wild quadratic extension of different
exponent two: the third sheet of the original trigonal cover is
unramified, and the ramified local degree-two extension is already
Galois. Thus $Z\to\mathbf P^1$ has total different degree two,
and $Z$ is rational. Choose a coordinate $z$ on $Z$ in which the
nontrivial involution is $z\mapsto z+1$, whose only fixed point is
infinity.

The cyclic cubic cover $Y\to Z$ is unramified above infinity and
ramified at six distinct finite points, in three pairs interchanged
by $z\mapsto z+1$. The two exponents in each pair are opposite,
since this involution acts on $C_3$ by inversion. Choosing the
positive exponent in each pair gives a monic separable cubic
\[
 F(z)=z^3+a z^2+b z+c,\qquad
 F^\sigma(z)=F(z+1),\qquad
                    y^3=F(z)/F^\sigma(z).
\]
Here $F$ and $F^\sigma$ have disjoint roots. Multiplying a Kummer
coordinate by a cube root of unity makes the chosen involution of
$Y$ exactly
\[
                         \sigma(z,y)=(z+1,y^{-1}).
\]
The quotient $Y/\langle\sigma\rangle$ is the original curve $C$,
up to the choice of a conjugate order-two subgroup of $S_3$.

The curve $Y$ has genus four. A basis of its regular differentials is
\[
 \omega_i=\frac{z^i\,dz}{yF^\sigma},\qquad
 \omega'_i=\frac{z^i\,dz}{y^2F^\sigma},\qquad i=0,1.
\]
To verify regularity directly, at a root of $F$ the respective orders
are one and zero, and at a root of $F^\sigma$ they are zero and one.
At each of the three points over infinity the orders are $1-i$.
The four forms are independent by their two distinct $C_3$
characters and the independence of $1,z$.
The involution sends
\[
 \omega_0\longleftrightarrow\omega'_0,\qquad
 \sigma\omega_1=\omega'_1+\omega'_0,\qquad
 \sigma\omega'_1=\omega_1+\omega_0.
\]
Its invariant subspace is thus two-dimensional, with basis
\[
       \Omega_0=\omega_0+\omega'_0,\qquad
       \Omega_1=\omega_1+\omega'_1+\omega'_0.
\]
Pullback along the separable map $Y\to C$ injects the two-dimensional
space $H^0(C,\omega_C)$ into this invariant subspace, hence identifies
it with the whole subspace. Cartier commutes with this pullback.
This also follows directly from its defining square-semilinearity
and the identities $C(dh)=0$ and $C(h\,dh)=dh$: a separating
parameter on $C$ remains a separating parameter in the separable
function-field extension.

Set $d=a+b+1$. Then
\[
 F^\sigma=z^3+(a+1)z^2+(b+1)z+(c+d).
\]
The four Cartier computations are
\begin{align*}
 C\omega_0&=\sqrt b\,\omega'_0+\omega'_1,&
 C\omega_1&=\sqrt c\,\omega'_0+\sqrt a\,\omega'_1,\\
 C\omega'_0&=\sqrt{b+1}\,\omega_0+\omega_1,&
 C\omega'_1&=\sqrt{c+d}\,\omega_0+\sqrt{a+1}\,\omega_1.
\end{align*}
In the invariant basis $(\Omega_0,\Omega_1)$ the Cartier matrix is
\[
           H=\begin{pmatrix}
                \sqrt{b+1}&\sqrt{a+c}\\
                1&\sqrt a
              \end{pmatrix},\qquad
                         (\det H)^2=ab+c.
\]
The last scalar is nonzero. Indeed $F'(z)=z^2+b$ has its unique
geometric root $\sqrt b$, and $F(\sqrt b)=ab+c$; its vanishing
would contradict separability of $F$. Thus Cartier on $C$ is
bijective in the $S_3$ case as well.
\end{proof}

\begin{corollary}[No triple three-torsion packet on a nonordinary theta curve]
\label{cor:nonordinary-genus-two-no-three-torsion-triple-candidate}
\textup{[Proved under the stated hypotheses; complete independent audits passed.]}
Let $C$ be nonordinary of genus two in characteristic two. There
are no three distinct points $P,Q,R\in C$ for which
\[
                         3P\sim3Q\sim3R.
\]
Consequently, for any subgroup $K\subset\operatorname{Jac}(C)[3]$,
each coset of $K$ meets an Abel--Jacobi theta curve in at most two
geometric points. This conclusion does not require isotropy of $K$.
\end{corollary}
\begin{proof}
Riemann--Roch gives $h^0(C,\mathcal O_C(3Q))=2$. The three stated
effective divisors define three distinct members of this pencil.
It is basepoint-free because two of these divisors have disjoint
supports. It therefore gives a degree-three map totally ramified
at all three points, contradicting the lemma. Points in a common
$K$-coset have differences annihilated by three, which gives the
same linear equivalences.
\end{proof}

\begin{corollary}[Every three-torsion theta norm is nodal]
\label{cor:nonordinary-theta-three-torsion-nodal-norm-candidate}
\textup{[Proved under the stated hypotheses; complete independent audits passed.]}
Let $B$ be a nonordinary indecomposable principally polarized
abelian surface in characteristic two, with a smooth symmetric
theta curve $C$, and let $K\subset B[3]$. The norm of $C$ under
$B\to B/K$ is integral and nodal, its normalization is $C$,
and it has exactly $|K|-1$ nodes. For each unordered nonzero type
$\{g,-g\}$ there are exactly two nodes.
\end{corollary}
\begin{proof}
Fix a Weierstrass point $W$ to identify $C$ with its Abel--Jacobi
image. The restriction of invariant differentials on $B$ to $C$
is an isomorphism onto $H^0(C,\omega_C)$, so equality of its
translated tangent lines is equivalent to equality of canonical
images. Thus a tangency between $C$ and $C+g$, with $g\ne0$,
would give $g=2(P-W)$ for a point $P$ with $P\ne\iota P$.
Since $3g=0$, the class $3(P-W)$ is geometric two-torsion.

If the two-rank is zero, this class is zero, so $3P\sim3W$.
But $|3W|=W+|K_C|$ by Riemann--Roch and $2W\sim K_C$;
thus $P=W$, a contradiction. If the two-rank is one, the
characteristic-two normal form in
Lemma~\ref{lem:p-rank-one-genus-two-polar-stable-candidate}
has exactly two Weierstrass points $W,W'$. Their nonzero
two-torsion difference generates the geometric two-torsion.
The remaining possibility is $3(P-W)=W'-W$, hence
$3P\sim K_C+W'$. Again $|K_C+W'|=W'+|K_C|$ by
Riemann--Roch, so $P=W'$, which is fixed by $\iota$ and gives
$g=0$. Thus every nonzero kernel translate intersects $C$
transversely, in exactly two points since $C^2=2$.

The preceding corollary excludes all three-point fibres of the
norm map. It is generically one-to-one: a nonzero translation
cannot stabilize the principal divisor, since the kernel of its
polarization is trivial. The quotient is \etale, so the local
image of the smooth curve has precisely ordinary nodes, each
arising from one pair of transverse branches. The two points of
$C\cap(C+g)$ give two such nodes for each type $\{g,-g\}$;
no identifications among these nodes remain because a fibre has
at most two points. The total is therefore $|K|-1$.
\end{proof}

\begin{remark}[Exact polynomial verification]
\label{rem:three-total-ramification-Cartier-certificate-candidate}
\textup{[Proved under the stated hypotheses; complete independent audits passed.]}
The executed program
\path{agents/monodromy/code/genus_two_three_total_ramification_checks.py},
with full output
\path{agents/monodromy/outputs/genus_two_three_total_ramification_checks.json},
reconstructs all four Cartier columns from odd polynomial coefficients,
checks the involution and its two-dimensional fixed basis, verifies
the restricted matrix, and computes the cubic separability resultant
$(ab+c)^2$. The author run at 19:25:11.722810 UTC on 19 November
2026 passed in 0.06777819900162285 seconds. These are universal
polynomial identities; the Galois-closure and ramification arguments
are proved above and are not inferred from finite examples.
\end{remark}

\begin{remark}[Scope]
\label{rem:nonordinary-three-torsion-triple-scope-candidate}
\textup{[Proved under the stated hypotheses; complete independent audits passed.]}
This establishes nodality of individual theta norms for kernels killed
by three, not for general odd kernels. It does not by itself prove
independence of node permutations or the gluing-translation input of the three-theta
conditional monodromy theorem. The calculation is intrinsic to
characteristic two; no characteristic-zero nodality assertion is
being specialized.
\end{remark}
